% ID: MAT-GEU-POL-001
% Titolo: Ottagono e volume di un cubo
% Disciplina: Matematica
% Area: geometria_euclidea
% Argomento: poligoni_regolari
% Classe: Prima secondaria di secondo grado
% Difficolta: 2
% Tipo: dimostrazione
% Tag: matematica, geometria_euclidea, poligoni_regolari
% Fonte: appunti_matematica_anonimizzati
% Provenienza: materiale fornito dal docente; testo anonimizzato; soluzione redatta per Exergo

\begin{esercizio}
Un ottagono regolare ha perimetro $2^{18}\,\mathrm{cm}$. Verifica che un cubo con spigolo uguale al lato dell'ottagono ha volume $2^{45}\,\mathrm{cm}^3$.
\end{esercizio}

\begin{soluzione}
Il lato misura $2^{18}/8=2^{15}\,\mathrm{cm}$. Lo spigolo del cubo ha questa lunghezza, quindi $V=(2^{15})^3=2^{45}\,\mathrm{cm}^3$. Il volume si esprime in centimetri cubi, non in centimetri.
\end{soluzione}
